% Reproduced inherited proof; finding P-F24.
\subsection{The chain contributes a smooth fixed component}

\begin{proposition}[P-F24: full chain stabilizer and local fixed dimension]
\label{prop:chainsmoothcomponent}
Let $n\ge4$, $M=2^{n-1}$, $g=\gcd(n,3)$, and let $F_n$ be the
chain cubic of P-F19. Its full projective linear automorphism group
is the diagonal cyclic group of order $M$. Its stabilizer as an
actual cubic form under $SL_n$ is exactly the scalar subgroup
$\boldsymbol\mu_g$, which acts trivially on all cubic forms.

Let $X=\operatorname{Spec}T^{(n)}$, with the stretch grading of
P-F19, and let $p$ be the image of $F_n$ in this quotient. Then
$X$ is smooth at $p$. The fixed scheme for the grading action of
$\boldsymbol\mu_M$ is smooth of dimension one at $p$.
The closure of the grading orbit of $p$ is consequently an
irreducible one-dimensional component of this fixed scheme,
and its nonzero grading orbit consists of smooth points of $X$.
If $3\nmid n$, this component has coordinate ring
$\mathbb C[z]$ with $\deg z=M$; its regraded leading Hilbert
coefficient $e_Y$ in P-F23 is therefore one.
This does not classify the other components of the fixed scheme.
\end{proposition}

\begin{proof}
We first determine the full projective automorphism group; computing
only its diagonal subgroup would not suffice. Half the Hessian matrix
of $F_n$ is the symmetric tridiagonal matrix with diagonal
\[
 (x_2,x_3,\ldots,x_n,3x_n)
\]
and successive off-diagonal entries $x_1,\ldots,x_{n-1}$.
We claim that its projective rank-at-most-two locus consists exactly of
\[
 [e_1],\qquad
 [a e_{n-1}+b e_n]\quad
 \text{with }a=0\text{ or }a^2=3b^2.
 \tag{P-F24a}
\]
The latter are three distinct points on a line not containing $[e_1]$.

Here is a full rank check. A tridiagonal block of length $l$ with
all its $l-1$ off-diagonal entries nonzero has kernel dimension at
most one: its first equation determines the second coordinate from
the first, and the subsequent equations determine the rest in order.
The zero off-diagonal entries divide the whole matrix into such
blocks. Its rank is therefore at least the number of nonzero
entries among $x_1,\ldots,x_{n-1}$. At rank at most two, at most
two of these coordinates are nonzero.

If precisely two are nonzero, say $x_i,x_j$ with $i<j$, first
suppose $j\ge i+2$. The block on indices $i,i+1$ has determinant
$-x_i^2$: its first diagonal entry $x_{i+1}$ is zero. It has
rank two. The separate block containing the other nonzero
off-diagonal has rank at least one, a contradiction. If
$j=i+1$ and $i\ge2$, the singleton on index $i-1$ has nonzero
diagonal $x_i$ and is separate from the length-three block on
$i,i+1,i+2$. Their ranks are at least one and two. If $i=1$,
then for $n\ge5$ the block on $1,2,3$ has determinant
$-x_2^3\ne0$. For $n=4$ this remains true when $x_4=0$;
when $x_4\ne0$, the singleton on index four has rank one and
the length-three block has rank at least two. This excludes
every case with two nonzero off-diagonal entries.

If the only nonzero one is $x_j$ with $2\le j\le n-2$,
the block on $j,j+1$ has determinant $-x_j^2$, and the
singleton on $j-1$ has diagonal $x_j$, again giving rank
at least three. If $j=1$, its length-two block already has
rank two; $x_n\ne0$ would contribute nonzero diagonal entries
on the disjoint last two indices. Thus only $[e_1]$ remains.
If $j=n-1$, put $a=x_{n-1}$ and $b=x_n$. The only nonzero
part is the singleton $a$ on index $n-2$ and the block
\[
 \begin{pmatrix}b&a\\a&3b\end{pmatrix}.
\]
Its total rank is two exactly when $a^2=3b^2$.
Finally, if there are no nonzero off-diagonal entries, the
projective point is $[e_n]$ and the rank is two. This proves
(P-F24a).

Under a linear change of coordinates the Hessian transforms by
congruence, so any projective linear automorphism of $F_n$
preserves (P-F24a). The three collinear points form its unique
collinear triple. Hence its remaining point $[e_1]$ is fixed.
Write a representing linear transformation in coordinate forms
$\ell_1,\ldots,\ell_n$, with
\[
 F_n(\ell_1,\ldots,\ell_n)=cF_n(x_1,\ldots,x_n),\qquad c\ne0.
\]
Fixing $[e_1]$ means $\ell_1=ax_1+L(x_2,\ldots,x_n)$,
$a\ne0$, while the other $\ell_i$ do not involve $x_1$.
Comparing the coefficients of $x_1^2$ gives
$a^2\ell_2=cx_2$, so $\ell_2=bx_2$ with $b\ne0$.
The coefficient of $x_1$ is $2aL\ell_2$ and must vanish.
Thus $L=0$. Restriction to $x_1=0$ leaves an automorphism
of the chain cubic in $n-1$ variables, whose first coordinate
form is a scalar multiple of its first variable.

This proves by induction that all the coordinate forms are
diagonal once the base case $n=4$ is checked. For that case
write the remaining ternary cubic as
\[
 G(y,z,w)=y^2z+z^2w+w^3.
\]
Its transformation has coordinate forms $by$, $uy+A(z,w)$,
and $vy+B(z,w)$, where $A,B$ are independent linear forms.
The coefficient of $y$ in its transformed cubic is
\[
 vA^2+2uAB+3vB^2.
\]
It is zero identically, so independence of $A,B$ implies
$u=v=0$. Comparison of the coefficient of $y^2$ now gives
$b^2A=cz$, so $A$ is a nonzero multiple of $z$.
Write $B=\alpha z+\beta w$, with $\beta\ne0$. The
coefficient of $zw^2$ is $3\alpha\beta^2$; hence
$\alpha=0$. This is the required diagonal form and completes
the induction.

For diagonal projective transformations choose a scalar
representative with last entry one. The last monomial then
forces the common multiplier to be one, and the other
monomials give $a_{i+1}=a_i^{-2}$. The last-entry condition is
$a_1^{(-2)^{n-1}}=1$. Thus the full projective group is
cyclic of order $M$. For linear transformations preserving
$F_n$ exactly, rather than its projective class, the same
recurrence instead ends with $a_n^3=1$. They form the
cyclic group of order $3M$, with entries
\[
 a_i=a_1^{(-2)^{i-1}},\qquad a_1^{3M}=1.
\]
Their determinant is $a_1^S$, where
\[
 S=1-2+4-\cdots+(-2)^{n-1}
   =\frac{1-(-2)^n}{3}.
\]
The integer $S$ is odd and satisfies $S\equiv n\pmod3$.
Consequently $\gcd(S,3M)=g$. The determinant-one kernel
has order $g$; when $g=3$, its entries are all the same
cube root of unity, because $(-2)^i\equiv1\pmod3$.
This proves the asserted $SL_n$ stabilizer.

By P-F19 the chain hypersurface is smooth. The standard
stability theorem for smooth hypersurfaces of degree at
least three implies that its nonzero affine $SL_n$ orbit
is closed. An exact source is Huybrechts, \emph{The geometry
of cubic hypersurfaces}, Chapter~3, Corollary~1.6, printed
pp.~149--150, together with Exercise~1.4 on p.~149 for the
affine formulation. The cited degree and dimension conditions
hold here, since $n\ge4$. Passing to $SL_n/\boldsymbol\mu_g$
does not change the quotient and makes the stabilizer trivial.
Luna's \emph{Slices \'etales}, Section~III.1, the unnumbered
slice theorem on p.~97 and Remark~1$^\circ$ on pp.~97--98, identify an
\'etale neighborhood of the quotient point with a neighborhood
in the quotient of a smooth normal slice by its stabilizer.
The stabilizer here is trivial, so the slice and the quotient
are smooth. Thus $X$ is smooth at $p$.

It remains to determine the local fixed dimension; smoothness
alone does not give it. Let $\zeta$ be primitive of order $M$,
and put
\[
 A=\operatorname{diag}
       (\zeta,\zeta^{-2},\ldots,\zeta^{(-2)^{n-1}}).
\]
This transformation fixes $F_n$. Its determinant is
$\zeta^S$ with $S$ odd. On degree-$s$ elements of the
canonical algebra $T^{(n)}$, it acts by the character
$\det(A)^{-3s/g}$, up to replacing $A$ by its inverse
according to the coordinate-action convention. Since
$3/g$ and $S$ are odd, this character generates exactly
the grading $\boldsymbol\mu_M$ action.

The only degree-three monomials fixed by $A$ are the $n$
monomials of $F_n$. To see this, the nonzero residue weights
$(-2)^i\pmod M$, $0\le i\le n-2$, have distinct
2-adic valuations, and the remaining weight is zero.
In a sum of three such weights that is zero modulo $M$,
the least valuation cannot occur only once. Its index
therefore occurs twice. If it is $i<n-1$, the third
weight must be $-2(-2)^i=(-2)^{i+1}$, giving exactly
a chain monomial. At the last nonzero weight, twice that
weight is zero and the third is the zero weight. Three
zero weights give $x_n^3$; three equal nonzero weights
never sum to zero because three is odd. This exhausts
the possibilities. Thus the fixed subspace of the cubic
representation has dimension $n$.

All eigenvalues of $A$ on the underlying $n$-space are
distinct, so its fixed subspace on $\mathfrak{sl}_n$
under conjugation consists of the traceless diagonal
matrices and has dimension $n-1$. The map
$\mathfrak{sl}_n\to T_{F_n}(SL_nF_n)$ is an isomorphism,
because the stabilizer is finite. Luna's smooth quotient
description, or the corresponding differential of the
local principal bundle, gives the exact tangent sequence
\[
 0\longrightarrow\mathfrak{sl}_n F_n
 \longrightarrow W^*
 \longrightarrow T_pX\longrightarrow0.
\]
It is $A$-equivariant. Taking fixed vectors is exact for
this finite-order action in characteristic zero. Hence
\[
 \dim(T_pX)^{\boldsymbol\mu_M}=n-(n-1)=1.
 \tag{P-F24b}
\]
The fixed scheme of a finite-order automorphism of a
smooth complex variety is smooth near a fixed point:
averaging local parameters into eigenspaces and passing
to the completed local ring identifies its fixed ideal
with the parameters of nontrivial character. This is
also the local argument used in P-F23. It follows from
(P-F24b) that the fixed scheme is smooth of dimension one
at $p$.

The positive-degree discriminant is nonzero at $p$, so
the grading orbit of $p$ has dimension one. All of it
is fixed by $\boldsymbol\mu_M$, since the grading group
is abelian. Its irreducible closure is contained in a
component of the fixed scheme through $p$. That component
has dimension one locally at its smooth point $p$, and
therefore has dimension one globally. The two irreducible
closed sets agree. Grading translations preserve smoothness
of $X$, so the entire nonzero orbit consists of smooth
points.

For the last assertion assume $3\nmid n$, so $g=1$.
Restriction to the grading orbit is the homomorphism
\[
 T^{(n)}\longrightarrow\mathbb C[u,u^{-1}],\qquad
 f\in T_s\longmapsto f(p)u^s.
\]
All exponents are nonnegative, and P-F19 shows that a
nonzero image has $M\mid s$. The discriminant has stretch
degree exactly $M$ and is nonzero at $p$. The image therefore
contains a nonzero multiple of $u^M$ and is exactly
$\mathbb C[u^M]$. This image is the coordinate ring of
the orbit closure, proving the claim. Its Hilbert series is
$1/(1-t^M)$, so $e_Y=1$. This completes the proof.
\end{proof}

\paragraph{The geometric gate isolated at this step.}
The diagonal subgroup itself is a known computation:
Zheng, \emph{On abelian automorphism groups of hypersurfaces},
arXiv:\texttt{2004.09008v3}, Lemma~3.2, printed p.~5,
computes the type-$T_n$ diagonal group. The Hessian argument
above verifies the full-group assertion needed here.
No historical priority is claimed. The local result establishes
the generic-smoothness hypothesis of P-F23 for the chain component.
What remains is to control all the other components of
$X^{\boldsymbol\mu_M}$, including those represented by
singular polystable cubics, sufficiently to prevent their
leading contributions from cancelling the chain contribution.
Uniqueness of the maximal-dimensional component would suffice.
This local proposition alone does not assert it; P-F30--P-F31
below establish it by a complete support classification and
thereby close the global noncancellation gate.

Primary sources for the cited geometric theorems are
\url{https://www.math.uni-bonn.de/people/huybrech/Notes.pdf}
and \url{https://www.numdam.org/item/10.24033/msmf.110.pdf}.
